\section{関連文献}
\label{sec:literature}

本稿は、それぞれ十分に発展しているが異なるmarginを重視する3つの文献群に関連する。すなわち、標準をめぐる政府の戦略的政策、私的互換性と企業適応、地域協力と多国間協力である。本稿の貢献は、新しい標準制度、互換性技術、coalition solution conceptを導入することではない。正式な標準化の取り決めの後に戦略的な私的互換性の意思決定が続き、政府がその結果として得られる継続利得を評価するときに生じるfeedbackにある。本節では、この選択的侵食メカニズムを最も近い先行研究との関係で位置付ける。

\subsection{標準、調和、政府の戦略的選択}
\label{subsec:lit-government-standards}

一つの文献群は、国際市場における戦略的政策手段として標準を研究する。\citet{gandalShy2001} は本稿の環境に特に近い。彼らは、conversion costとnetwork effectを伴う3か国設定で外国標準に対する政府のrecognitionを分析し、国がstandardization unionを形成することを認める。conversion costが十分に大きいとき、2国は第3国の標準を排除するunionから利益を得ることができる。したがって彼らの分析は、compatibility costと政府のrecognition policyがexclusive regional standards arrangementを生み得ることを確立する。しかし彼らのモデルでは、nonrecognitionに直面する外国企業は、仕向地市場で販売するためにconversion costを負担しなければならない。conversion responseはrecognition regimeに結び付いており、正式coalitionが固定された後に、標準固有の固定費用を伴う独立したadoption gameを通じて選ばれるわけではない。本稿は後者のmarginを内生化し、その誘導されたcontinuation stateがcoalition stabilityへどのようにfeedbackするかを問う。

他の研究は、企業、政府、国際標準の間の異なる戦略的緊張を明らかにしている。\citet{barrettYang2001} は、redesign costとnetwork effectによって、既存企業が国際製品標準からdeviateすることが合理的となり、host governmentがその標準をenforceする誘因も弱まり得ることを示す。\citet{klimenko2009} は、network externalityの下でtechnical compatibilityをめぐる政府の戦略的政策を研究し、compatibility policyに関する国際協定を評価する。密接に関連する2か国設定で、\citet{klimenkoInteroperability2009} は、政府のinteroperability standard、trade tax、外国企業による費用のかかるcompatibility-enhancing investmentを分析する。これらの研究はcompatibilityが企業行動とpublic policy incentiveの双方に影響し得ることを明らかにする。本稿の問いはより狭く、正式な地域coalitionがすでに存在する後に企業が追加標準を私的に選択する応答と、その応答がcoalition stabilityへ与える影響に関する。

関連するtrade-agreement literatureは、政府が戦略的標準と国内規制をどのように規律できるかを研究する。\citet{bagwellStaiger2001} はtrade agreementのmarket-access logicを通じて国内政策に対する国際的規律を分析する。\citet{costinot2008} は、horizontal compatibility standardを含むoligopoly environmentで、product standardに関するnational treatmentとmutual recognitionを比較する。\citet{geng2019} も、consumption externalityに対する国の選好が異なる場合にnational treatmentとmutual recognitionを比較する。より最近では、\citet{grossmanMcCalmanStaiger2021} が、企業が製品variantを各ローカル市場にtailorし、その違いが固定費用を生む場合のregulatory convergenceを研究する。これらの研究は国際的regulatory ruleの設計または効率性を中心とする。これに対し本稿は、単純なformal standards partitionのmenuを固定し、企業の内生的な技術的応答がその分割に付随するnational-welfare payoffをどのように変えるかを研究する。

regional-versus-multilateral standards marginについて最も近い文献は、\citet{takaradaEtAl2020} と \citet{kawabataTakarada2021} である。\citet{takaradaEtAl2020} はnational、regional、multilateral standardを備えた3か国Cournot modelを構築し、coreを用いてどのstandards regimeが維持され得るかを特徴付ける。また、異なる標準の下で生産をsetupする固定費用も考える。\citet{kawabataTakarada2021} はtariffを加え、deep free trade agreementとcustoms unionを区別し、regional harmonizationがmultilateral harmonizationのbuilding blockかstumbling blockかを問う。これらの研究は、regional standardがmultilateral cooperationと非自明に相互作用し得ることを確立している。本稿との違いは、固定費用そのものの存在ではない。\citet{takaradaEtAl2020} のfixed-cost extensionでは、企業はすべての市場へ供給すると仮定され、支配的な政府標準によって複数標準の下での生産が必要になる場合にsetup costを負う。本稿では、追加標準のsupportそのものが企業の戦略的意思決定である。この違いにより、multi-standard productionがformal regimeによって機械的に課される場合には存在しないcontinuation-state comparisonが生じる。

\subsection{私的互換性と企業適応}
\label{subsec:lit-private-compatibility}

第2の文献群は、私的な戦略的選択としてcompatibilityを研究する。\citet{katzShapiro1985} のclassic network-competition frameworkは、消費者がnetwork sizeを評価する場合、compatibilityが需要と競争をどのように変えるかを示す。\citet{besenFarrell1994} は、戦略環境によっては企業自身がcompatible competitionまたはincompatible competitionを選好し得ることを強調する。最も直接的には、\citet{farrellSaloner1992} が、本来互換でないtechnologiesのinteroperationを可能にするconverterを分析する。彼らの分析は、private deviceがstandardまたはinterfaceによって生じたincompatibilityを部分的に克服できることを確立する。したがって、非互換性に対するprivate workaroundそれ自体は本稿の新規性ではない。

正式な互換性とnetwork benefitへアクセスする代替手段との区別は、\citet{doganogluWright2006} においても中心的である。彼らはmultihomingがcompatibilityを不完全に代替し得ること、またこのtechnological substitutionがcompetitionとcompatibilityのsocial desirabilityを変え得ることを示す。これは本稿が「technological substitute」という表現を用いる際の重要な先行研究である。本稿のpolitical-complement interpretationは対象が異なる。それは、私的互換性がregional exclusion rentを除去したときに生じるgovernment-incentive cross-differenceだけを指す。compatibilityとformal harmonizationが一般的なgame-theoretic senseでcomplementsであるという主張ではない。

貿易モデルにも、標準をまたぐ費用のかかる企業適応についての先行研究がある。\citet{fischerSerra2000} は、外国producerが複数のstandard levelで生産すると固定setup costを負うことを認め、その費用がprotectionist standards policyとどのように相互作用するかを示す。\citet{schmidtSteingress2022} は、新しいcross-country standards dataとheterogeneous-firm modelを組み合わせ、企業がsunk adoption costを負ってharmonized standardを採用するかを選ぶモデルを構築する。harmonizationはcountry-specific adaptation costを削減し、tradeを拡大する。これらの研究は、multi-standard production costとendogenous firm adoptionの双方が経済的に重要なmarginであることを確立している。本稿のメカニズムは、そのどちらかを初めて導入することに依存しない。重要なのは、追加標準への投資の収益性が現在のformal partitionに依存し、結果として生じる非対称なadoption patternが、後続のcoalition-stability comparisonで政府が用いるcontinuation payoffを変えることである。

regulatory diversityに関する最近の研究も、同じclaim disciplineを補強する。\citet{parentiVannoorenberghe2025} は、企業がcountry-specific standardに製品をtailorすると費用を負うimperfectly competitive trade modelを構築し、regulatory cooperationがblocを形成し得ることを示す。\citet{maggiMrazova2024} はregulatory diversityの固定費用を研究し、regulatory agreementがなくてもharmonizationが自発的に生じ得ることを示す。これらの貢献は、異質なregulatory environmentへ供給する企業レベルの費用それ自体がharmonization incentiveを変え得ることを示している。本稿の結果は構造が異なる。域外国企業がregional bloc standardを選択的に採用でき、その結果、coalition加盟企業による域外国市場での相互採用を誘発することなく、加盟国市場で生じるcoalition-specific benefitを消滅させるのである。

\subsection{地域主義、regulatory bloc、coalition formation}
\label{subsec:lit-coalitions}

第3の関連はcoalition formationとregionalism-versus-multilateralism questionである。standards literatureにはすでに明示的なcoalition comparisonが存在する。前述のとおり、\citet{takaradaEtAl2020} はcore conceptを用いてnational、regional、multilateral standards regimeを比較し、\citet{kawabataTakarada2021} はdeep regional agreementがmultilateral agreementを促進または阻害する条件を特徴付ける。\citet{parentiVannoorenberghe2025} は、regulatory preferenceが十分に分散しているとregulatory blocが維持され得ることを示す。したがって、本稿はregional standard、standards coalition、building-block/stumbling-block questionが新しい対象であるとは主張しない。

本稿は新しいcoalition solution conceptを提案するものでもない。partition-based coalition formationには長い理論的伝統があり、\citet{hartKurz1983} はendogenous coalition structureのcanonical analysisを与える。本稿では、deliberately simpleなexclusive-membership partition gameとstrict blockingを用いる。関心のあるeconomic objectは、むしろ先行するindustrial-organization stageがcoalition gameへ供給するcontinuation payoffである。formal membershipはbaseline compatibility structureを決めるが、その後企業がadditional standardをサポートするかを選ぶ。そのためcoalitionは、product-market behaviorとprivate technological adaptationがformal arrangementに反応した後で評価される。このtimingにより、政府段階では従来型のblocking logicを維持しつつ、coalition payoffをpost-formation firm conductに対して内生化できる。

interoperabilityとcoalition formationはstandards-trade literatureの外でも一緒に現れる。\citet{guoLiuNault2024} は、加盟者がより高いinteroperabilityを享受し、その後coalitionがresource investmentを選ぶcritical IT infrastructureのcoalition formationを研究する。彼らのモデルではinteroperabilityがcoalition membershipの経済的に重要な特徴だが、formal coalition membershipを固定したままeffective compatibilityを変えるfirm-level post-formation bypassは存在しない。この比較は本稿のclaimをさらに限定する。interoperabilityとcoalition formationそれ自体には先行研究があり、本稿の対象は既存のstandards boundaryに対するstrategic private circumventionが生み出すcontinuation-payoff feedbackである。

この区別は、政府協定自体が適用されるstandardまたはrecognition ruleを決めるregulatory-agreement modelからも本稿を分ける。\citet{costinot2008}、\citet{geng2019}、\citet{grossmanMcCalmanStaiger2021} では、制度比較の対象はnational treatment、mutual recognition、coordinated regulationなどのruleである。本モデルでは、private adoptionはformal agreementを修正しない。同じregional partitionでも、域外国企業がbloc standardを有利にサポートできるかどうかに応じて異なるeffective compatibility patternを生み得る。したがってinstitutional menuと各partitionのformal membershipが変わらなくても、coalition stabilityは変わり得る。

\subsection{選択的侵食メカニズムの位置付け}
\label{subsec:lit-positioning}

以上の最も近い先行研究は、周辺を構成する要素を提供している。government standards policyはstandardization unionとstrategic exclusionを生み得る \citep{gandalShy2001}。government interoperability policyは費用のかかるfirm compatibility investmentと相互作用し得る \citep{klimenkoInteroperability2009}。regional standardとmultilateral standardは、fixed multi-standard production costがある場合を含め、coalition stabilityを用いて比較されている \citep{takaradaEtAl2020}。private compatibility technologyはincompatibilityを克服し得る \citep{farrellSaloner1992,doganogluWright2006}。また、企業はharmonized standardを内生的に採用したり、regulatory diversityの費用を負ったりし得る \citep{schmidtSteingress2022,maggiMrazova2024}。本稿の貢献は、これらのmargin間に生じる特定のfeedbackである。

地域標準化連合は当初、coalition加盟国市場で域外国企業を排除することに伴うnational-welfare componentを加盟国にもたらす一方、coalition加盟国は域外国市場で不利益を負う。次に域外国企業はbloc standardを私的にサポートすることを選べる。その意思決定はformal coalition membershipを変えず、coalition加盟企業に域外国標準をサポートさせることもなく、加盟国市場でのproduct-market competitionを変える。そのためcontinuation payoffは非対称に変化する。加盟国市場の排除構成要素は消えるが、相互的な域外国市場の不利益は残る。Sections~\ref{sec:selective-erosion} と \ref{sec:coalition-stability} では、この選択的侵食が政府の順位付けを逆転させ、特定されたopen parameter regionでstrict-blocking stable setをregional standards unionからinternational standardizationへ変え得ることを示す。

したがって本稿は、ingredientsの組合せではなくmechanismを中心に位置付ける。standards、fixed costs、network effects、governments、coalition formationを組み合わせたこと自体に新規性を主張しない。より狭いclaimは、strategic private compatibilityがformal exclusionのcontinuation valueを内生的に変え、その結果standards coalitionを維持する政治的誘因を変え得るというものである。これが、追加標準をサポートするかというindustrial-organization decisionを、政府がどのformal standards arrangementを維持する意思を持つかというinstitutional questionへ結び付けるchannelである。
